\documentclass{article}
\usepackage{../../math_template}

\author{Jonathan Kasongo}
\title{Noether Sheet 4}

\begin{document}
\maketitle
\centerline{
These were my solutions to the Noether Mentoring Scheme Sheet 4. Problem
statements are abridged.
}

\begin{problem}{Problem 1}
Solve the simultaneous equations
\[
\begin{aligned}
4x + 3y &= 496\\
\sqrt{5x} + \sqrt{5y} &= \sqrt{17x + 7y}.
\end{aligned}
\]
\end{problem}

\begin{solution}{Solution}
Define $\alpha = 496$. Square the second equation twice to end up with
$$
144x^2 - 52xy + 4y^2 = 0.
$$
Now substitute in $4x = \alpha - 3y$ to get
$$
9(\alpha - 3y)^2 - 13y(\alpha - 3y) + 4y^2 = 0
$$
which simplifies to
$$
124y^2 - 67\alpha y + 9\alpha^2 = 0
$$
which yields the solutions $y = 144, 124$, with corresponding values of
$x = 16, 31$, and both solutions can be checked to work. $\blacksquare$
\end{solution}
\newpage

\begin{problem}{Problem 2}
(a) Show algebraically that
$$
\binom{n+1}{r+1} = \binom{n}{r} + \binom{n}{r+1}.
$$
\vspace{0.1cm}

(b)
\vspace{-0.65cm}
\begin{enumerate}[label=\roman*.]
  \item How many ways are there to choose $r+1$ objects from a set of $n+1$.
  \item Suppose that one of the $n+1$ objects is special. How many ways can
  you choose the $r+1$ items including the special one?
  \item Hence give a combinatorial proof that $$
  \binom{n+1}{r+1} = \binom{n}{r} + \binom{n}{r+1}.
  $$
\end{enumerate}

(c) Show that
$$\binom{n}{0} + \binom{n}{1} + \cdots + \binom{n}{n} = 2^n.$$
\end{problem}

\begin{solution}{Solution}
(a) The RHS reads
\[
\begin{aligned}
\binom{n}{r} + \binom{n}{r+1} &= \frac{n!}{r!(n-r)!} +
\frac{n!}{(r+1)!(n-r-1)!}\\
&= \frac{n!}{r!} \left( \frac{1}{(n-r)!} + \frac{n-r}{(r+1)(n-r)!} \right)\\
&= \frac{n!}{r!} \left( \frac{(r+1) + (n-r)}{(r+1)(n-r)!} \right)\\
&= \frac{n!}{r!} \left( \frac{n+1}{(r+1)(n-r)!} \right)\\
&= \frac{(n+1)!}{(r+1)!(n-r)!}\\
&= \binom{n+1}{r+1}. \ \square
\end{aligned}
\]
\\

(b) i. There are $\binom{n+1}{r+1}$ ways to choose $r+1$ objects from a set
of $n+1$.
\vspace{0.2cm}

ii. We must pick the one special object, other than it, we need to pick
$r$ items from a set of $n$ items. There are $\binom{n}{r}$ ways to do this.
\vspace{0.2cm}

iii. From i and ii we know that there are $\binom{n+1}{r+1} - \binom{n}{r}$
ways to choose a set of $r+1$ objects from a set of $n+1$, excluding the
special one. Alternatively, to pick $r+1$ objects excluding the special one,
we can just remove the special object from the
set and pick $r+1$ objects from this new set. There $\binom{n}{r+1}$ ways
to do this. But this means that
$$
\binom{n+1}{r+1} - \binom{n}{r} = \binom{n}{r+1}
$$
which re-arranges to the desired result. $\square$
\\

(c) Note that the sum is exactly $(1+1)^n$ due to the binomial theorem
which is the same as $2^n$. $\blacksquare$
\end{solution}
\newpage

\begin{problem}{Problem 3}
Suppose that $l$ is a line, and that $\Gamma$ and $\Delta$ are circles,
tangent to $l$ at $C$ and $D$ respectively and also tangent to each other
at $T$. Let $E$ be the opposite end of the diameter from $\Delta$ from $D$.
\vspace{0.2cm}

Prove that $C, T$ and $E$ are collinear.
\end{problem}

\begin{solution}{Solution}
The following diagram gives an overview of our arguement.

\begin{center}
\begin{asy}
/* Geogebra to Asymptote conversion, documentation at artofproblemsolving.com/Wiki go to User:Azjps/geogebra */
import graph; size(220);
real labelscalefactor = 0.5; /* changes label-to-point distance */
pen dps = linewidth(0.7) + fontsize(10); defaultpen(dps); /* default pen style */
pen dotstyle = black; /* point style */
real xmin = -2.325167272386684, xmax = 10.129696578885018, ymin = -1.025734314251339, ymax = 8.815415888658772;  /* image dimensions */
pen ccqqqq = rgb(0.8,0.,0.); pen qqqqcc = rgb(0.,0.,0.8); pen qqwuqq = rgb(0.,0.39215686274509803,0.);
pair gamma  = (0.,2.), delta  = (5.656854249492381,4.), T = (1.8856180831641276,2.6666666666666665), D = (5.656854249492381,0.), C = (0.,0.), P = (2.8284271247461885,0.);

filldraw(arc(C,0.8085208357109156,414.7356103172453,450.)--(0.,0.)--cycle, ccqqqq + opacity(0.10000000149011612),  ccqqqq);
filldraw(arc((5.656854249492381,8.),0.8085208357109156,234.73561031724537,270.)--(5.656854249492381,8.)--cycle, ccqqqq + opacity(0.10000000149011612),  ccqqqq);
filldraw(arc(gamma ,0.8085208357109156,270.,379.4712206344907)--(0.,2.)--cycle, qqqqcc + opacity(0.10000000149011612),  qqqqcc);
filldraw(arc(delta ,0.8085208357109156,90.,199.4712206344907)--(5.656854249492381,4.)--cycle, qqqqcc + opacity(0.10000000149011612),  qqqqcc);
filldraw((1.346604192690184,2.476096478112732)--(1.537174381244118,1.9370825876387885)--(2.076188271718062,2.127652776192723)--T--cycle, qqwuqq + opacity(0.10000000149011612),  qqwuqq);
filldraw((2.0761882717180615,2.127652776192723)--(2.6152021621920056,2.318222964746657)--(2.4246319736380713,2.8572368552206004)--T--cycle, qqwuqq + opacity(0.10000000149011612),  qqwuqq);
/* draw figures */
draw(circle(gamma , 2.));
draw(circle(delta , 4.));
draw(gamma --C);
draw(delta --D);
draw((xmin, 0.*xmin + 0.)--(xmax, 0.*xmax + 0.)); /* line */
draw((xmin, -2.828427124746183*xmin + 8.)--(xmax, -2.828427124746183*xmax + 8.)); /* line */
draw(gamma --T);
draw(T--delta );
draw((5.656854249492381,8.)--delta );
draw(C--T,  linetype("4 4"));
draw(T--(5.656854249492381,8.),  linetype("4 4"));
/* dots and labels */
dot(gamma ,dotstyle);
label("$\gamma$", (0.07880339263036736,2.1849424514030202), NE * labelscalefactor);
dot(delta ,dotstyle);
label("$\delta$", (5.720482112924312,4.179293846156609), NE * labelscalefactor);
label("$\Gamma$", (-1.0171915179999849,3.35280588076323), NE * labelscalefactor);
label("$\Delta$", (3.654262199440861,7.054034595350973), NE * labelscalefactor);
dot(T, dotstyle);
label("$T$", (1.9653520092891705,2.813791990289287), NE * labelscalefactor);
dot(D, dotstyle);
label("$D$", (5.720482112924312,0.1366896676020362), NE * labelscalefactor);
dot(C, dotstyle);
label("$C$", (0.07880339263036736,0.1366896676020362), NE * labelscalefactor);
dot((5.656854249492381,8.), dotstyle);
label("$E$", (5.720482112924312,8.150029505981324), NE * labelscalefactor);
label("$l$", (-2.8678058752938584,0.154656797284501), NE * labelscalefactor);
label("$k$", (-0.8195530914928723,10.413887845971885), NE * labelscalefactor);
dot((5.656854249492385,0.), dotstyle);
//dot((3.142696805273544,7.11111111111112), dotstyle);
dot((1.8856180831641298,2.6666666666666643), dotstyle);
dot((1.8856180831641298,2.6666666666666643), dotstyle);
dot(P, dotstyle);
label("$P$", (2.8996427527773396,0.1366896676020362), NE * labelscalefactor);
clip((xmin,ymin)--(xmin,ymax)--(xmax,ymax)--(xmax,ymin)--cycle);
/* end of picture */
\end{asy}
\end{center}
\textit{Claim 1: $\gamma, T$ and $\delta$ are collinear.} \vspace{0.2cm}

\textit{Proof:} We note that because radii meet tangents at right angles
${\color{darkgreen}{\angle PT \gamma}} =
{\color{darkgreen}{\angle PT \delta}} = 90^\circ$. Now since
$\angle \gamma T \delta = 90 + 90 = 180^\circ$, we can conclude that
$\gamma, T$ and $\delta$ lie on a common line. $\square$ \\

\textit{Claim 2: ${\color{blue}{\angle C\gamma T}} =
{\color{blue}{\angle T\delta E}}$.}\vspace{0.2cm}

\textit{Proof:} Notice that since both $C\gamma$ and $D\delta$ are
perpendicular to the line $l$, we have $C\gamma \parallel D\delta$. Now due
to Alternate angle theorem
${\color{blue}{\angle C\gamma T}} = {\color{blue}{\angle T\delta E}}$.
$\square$
\\

Finally since $\triangle TC\gamma$ and $\triangle TE\delta$ are both
isoceles,
${\color{red}{\angle TC\gamma}} = {\color{red}{\angle TE\delta}}$.
Now since both $CT$ and
$TE$ make the the same angle with the with the line perpendicular to $l$, we
have $CT \parallel TE$ and since the two segments both pass through the
point $T$, it is the case that $C,T$ and $E$ are collinear. $\blacksquare$
\end{solution}
\newpage

\begin{problem}{Problem 4}
(a) What is the remainder when $1000^{1000}$ is divided by 89?
\vspace{0.2cm}

(b) What is the remainder when $1000^{1000}$ is divided by 87?
\end{problem}

\begin{solution}{Solution}
(a) We notice that 89 is a prime number. Since $1000 = 11\cdot 89 + 21$,
we wish to compute $21^{1000} \ (\text{mod } 89)$. We use Fermat's little
theorem,
$$
21^{1000} \equiv (21^{11})^{89}\cdot 21^{21} \equiv 21^{11 + 21}
\equiv 21^{32} \ (\text{mod } 89).
$$
Now since $21^2 = 121 \equiv 32 \ (\text{mod } 89)$ we have to compute
$32^{16} \ (\text{mod } 89)$. Now we can keep applying this same trick
to find
$$
32^{16} \equiv 45^8 \equiv 67^4 \equiv 39^2 \equiv 8 \ (\text{mod } 89). \
\square
$$
\phantom{.}\\
(b) We make use of Euler's theorem. Note that
$\varphi(87) = 87\left(1 - \frac{1}{3}\right)\left(1 - \frac{1}{29}\right)
= 56$. Since $3000 = 56\cdot 53 + 32$ we can write
$$
1000^{1000} = 10^{3000} = (10^{53})^{56} \cdot 10^{32}.
$$
Because 10 and 87 are coprime we can apply Euler's theorem to reduce the
last expression mod 87. It becomes
$$
(1)  10^{32} \ (\text{mod } 87).
$$
We can use the same trick as we used in part (a),
$$
10^{32} \equiv 13^{16} \equiv 82^8 \equiv 25^4 \equiv 16^2 \equiv 82 \
(\text{mod } 87). \ \blacksquare
$$
\end{solution}
\newpage

\begin{problem}{Problem 5}
Jose buys a lottery ticket for his $45$-th birthday, for which he has to
select six numbers between 1 and $59$.
To celebrate his age, he chooses the number $45$ as one of them. In his six
numbers, he uses all the
digits $0$ to $9$ once each, and the product of his numbers is a perfect
square.
\vspace{0.2cm}

What numbers did he pick?
\end{problem}

\begin{solution}{Solution}
We begin with a few claims.
\vspace{0.2cm}

\textit{Claim 1: The chosen number that uses the digit $7$ is either $7$ or $27$.}
\vspace{0.2cm}

\textit{Proof:} Let $x$ be the chosen number that uses the digit $7$. If $x$ has
two digits, then it's tens digit must be $1,2,3$ since Jose selects numbers from
$1$ to $59$. Note that $17$ and $37$ are prime numbers.
\begin{itemize}
\item If the number $17$ is chosen then, there must be another multiple of $17$
among the chosen numbers to make the product square. The only other multiples of
$17$ in the given range are $34, 51$. But $34$ can't be used since the digit $4$
has been taken, and $51$ can't be used since the digit $5$ has been taken.
\item If the number $37$ is chosen, then there must be another multiple of $37$
among the chosen numbers to make the product square. However since $37$ itself is
the only multiple of $37$ in the given range, the number $37$ can't be chosen.
\end{itemize}
Thus it must be the case that $x \in \{ 7, 27 \}$. $\square$
\\

\textit{Claim 2: The chosen number that uses the digit $0$ is either $10, 20$ or
$30$.}
\vspace{0.2cm}

\textit{Proof:} Since $45 = 3^2 \cdot 5$, one of the chosen numbers must have an
odd $5$-adic valuation, in order to ensure that the product is square, and in
particular such a number must be a multiple of $5$. Recall that multiples of $5$
can only end in the digits $0, 5$. Since the digit $5$ has been taken, such a
number must end in the digit $0$. Finally since all the numbers are between
$1$ and $59$ and the digits $4, 5$ have been taken, the only permissable
numbers that use the digit $0$ are $10, 20$ and $30$. $\square$
\\

\textit{Claim 3: The chosen number that uses the digit $7$, must be $27$.}
\vspace{0.2cm}

\textit{Proof:} Let $x$ be the chosen number that uses the digit $7$.
Thanks to claim 1, we need only demonstrate that $x \neq 7$. We proceed by way
of contradiction. Suppose that $x = 7$. Another one of the chosen numbers must
have an odd $7$-adic valuation in order to ensure that the product is square, and
in particular such a number must be a multiple of $7$. The only other multiples of
$7$ in the given range that don't use the digits $4,5,7$ are $21, 28$.
\vspace{0.2cm}

Suppose that the other multiple of $7$ is $21$. Since the digits $1,2$ are taken
another one of the chosen numbers must be $30$ as a consequence of claim 2.
The remaining digits are $6,8,9$, and we need to choose two more numbers. It is
easy to check that one of the numbers must have two digits and the other one must
have one digit. However, the two digit number is at least $68 > 59$, a
contradiction. \vspace{0.2cm}

Suppose that the other multiple of $7$ is $28$. Again due to claim 2, one of the
other numbers must be from the set $\{ 10, 30 \}$. If the number $10$ is chosen
then then the remaining digits are $3,6,9$. There are two numbers left to choose.
Since
$\nu_2 (45\cdot 7\cdot 28\cdot 10) = 3$, exactly one of the last two numbers must
have an
odd $2$-adic valuation. The only number that has an odd $2$-adic valuation is $6$,
so the number $6$ has to be chosen also. Thus the complete set of chosen numbers
is $\{ 45, 7, 28, 10, 6, 39 \}$. However since $39 = 3\cdot 13$ and this is the
only multiple of $13$ amongst the six numbers, the product isn't a perfect square,
a contradiction. $\square$ \\

So we must have chosen the numbers $45, 27$. Due to claim 2, we must've also
chosen one of the numbers $10, 30$. If we choose the number $10$ then the
remaining digits are $3,6,8,9$, and we have to choose three more numbers. It is
easy to check that one of the numbers must have two digits, and the other two
numbers must have one digit each. The two digit number must come from the set
$\{ 36, 38, 39 \}$. Now one can check that if we choose the number
$38 = 2 \cdot 19$ then the product won't be a perfect square since there is no
way to get an additional factor of $19$, similarly if we choose the number
$39 = 3 \cdot 13$ then the product won't be a perfect square since there is no way
to get an additional factor of $13$. Thus the two digit number is $36$, and the
complete set of chosen numbers is $\{ 45, 27, 10, 36, 8, 9 \}$. But since
$\nu_3(45\cdot 27\cdot 10\cdot 36\cdot 8\cdot 9) = 9$, the product isn't a
perfect square. Thus we cannot choose the number $10$.
\vspace{0.2cm}

Thus we must have chosen the numbers $45, 27, 30$. The remaining digits are
$1, 6, 8, 9$ and we need to choose three more numbers. Again we must
have 1 two digit number and 2 one digit numbers. The two digit number is one of
$16, 18, 19$, and from here it is easy to check case by case that the complete
set of chosen number is $\{ 45, 27, 30, 16, 8, 9 \}$, which has product
$45\cdot 27\cdot 30\cdot 16\cdot 8\cdot 9 = (6480)^2$. $\blacksquare$
\end{solution}
\newpage

\begin{problem}{Problem 6}
An \textit{equilateral} shape is one where all the sides have the same
length.\vspace{0.2cm}

(a) Construct an equilateral hexagon of side length 5 such that all the
vertices have integer coordinates.\vspace{0.2cm}

(b) Is it possible to construct an equilateral pentagon of side length 5
where all the vertices have integer coordinates?
\end{problem}

\begin{solution}{Solution}
(a) Consider the following construction.
\begin{center}
\begin{asy}
/* Geogebra to Asymptote conversion, documentation at artofproblemsolving.com/Wiki go to User:Azjps/geogebra */
import graph; size(220);
real labelscalefactor = 0.5; /* changes label-to-point distance */
pen dps = linewidth(0.7) + fontsize(10); defaultpen(dps); /* default pen style */
pen dotstyle = black; /* point style */
real xmin = -1.083722943722935, xmax = 14.167359307359302, ymin = -1.630068333869977, ymax = 7.385083181281503;  /* image dimensions */
pen xdxdff = rgb(0.49019607843137253,0.49019607843137253,1.); pen bubtba = rgb(0.7058823529411765,0.7019607843137254,0.7294117647058823); pen rcrcrf = rgb(0.10980392156862745,0.10980392156862745,0.12156862745098039);
pair A = (4.,0.), B = (9.,0.), C = (13.,3.), D = (9.,6.), F = (0.,3.);

filldraw((4.,6.)--D--C--B--A--F--cycle, xdxdff + opacity(0.10000000149011612),  xdxdff);
/* draw grid of horizontal/vertical lines */
pen gridstyle = linewidth(0.7) + bubtba; real gridx = 1., gridy = 1.; /* grid intervals */
for(real i = ceil(xmin/gridx)*gridx; i <= floor(xmax/gridx)*gridx; i += gridx)
draw((i,ymin)--(i,ymax), gridstyle);
for(real i = ceil(ymin/gridy)*gridy; i <= floor(ymax/gridy)*gridy; i += gridy)
draw((xmin,i)--(xmax,i), gridstyle);
/* end grid */

Label laxis; laxis.p = fontsize(10);
xaxis(xmin, xmax,defaultpen+rcrcrf, Ticks(laxis, Step = 1., Size = 2, NoZero),EndArrow(6), above = true);
yaxis(ymin, ymax,defaultpen+rcrcrf, Ticks(laxis, Step = 1., Size = 2, NoZero),EndArrow(6), above = true); /* draws axes; NoZero hides '0' label */
/* draw figures */
draw(F--(4.,6.));
draw((4.,6.)--D);
draw(D--C);
draw(C--B);
draw(B--A);
draw(A--F);
draw((4.,6.)--D,  xdxdff);
draw(D--C,  xdxdff);
draw(C--B,  xdxdff);
draw(B--A,  xdxdff);
draw(A--F,  xdxdff);
draw(F--(4.,6.),  xdxdff);
/* dots and labels */
dot(A,dotstyle);
label("$A$", (4.076450216450217,0.17079746699581105), NE * labelscalefactor);
dot(B,dotstyle);
label("$B$", (9.063463203463199,0.17079746699581105), NE * labelscalefactor);
dot(C,dotstyle);
label("$C$", (13.063463203463197,3.1664684626668), NE * labelscalefactor);
dot(D,dotstyle);
label("$D$", (9.063463203463199,6.179455475653806), NE * labelscalefactor);
dot((4.,6.),dotstyle);
label("$E$", (4.076450216450217,6.179455475653806), NE * labelscalefactor);
dot(F,dotstyle);
label("$F$", (0.07645021645022068,3.1664684626668), NE * labelscalefactor);
clip((xmin,ymin)--(xmin,ymax)--(xmax,ymax)--(xmax,ymin)--cycle);
/* end of picture */
\end{asy}
\end{center}
Each pair of adjacent points differs by the vector $\langle \pm 5,0\rangle$ or by the
vector $\langle \pm 4, \pm 3 \rangle$. Since both vectors have magnitude 5
this construction is valid. $\square$
\\

(b) Call an equilateral shape \textit{primitive} if all it's vertices lie
on lattice points. In what follows if we have $\pm$ in both the $x$ and $y$
component of a point/vector we assume that each sign is chosen
independently. \vspace{0.2cm}

\textit{Claim 1: On a primitive shape with side length 5, adjacent vertices
must differ by one of the vectors $\langle \pm 5,0\rangle$ or
$\langle \pm 3, \pm 4 \rangle$ and their permutations.}
\vspace{0.2cm}

\textit{Proof:} Let $P$ and $Q$ be adjacent vertices of a primitive shape.
Without loss of generality we can let one of the vertices
lie on the origin, because we can translate the primitive shape in such a
way that this occurs. Let $P = (0,0), Q = (a,b)$ for some
$a,b \in \mathbb{Z}$. Because the primitive shape has side length 5, we
have
$$
a^2 + b^2 = 5^2.
$$
Now one can test the squares $a^2 = 0,1,4$ and we find that $(a,b)$ is a
permutation of $(\pm 5, 0)$ or $(\pm 3, \pm 4)$, from which the desired
result follows. $\square$\\

Say that a vector is of type I if it is of the form
$\langle \pm 5, 0 \rangle$ or $\langle \pm 3, \pm 4\rangle$, say that a
vector is of type II if it is of the form $\langle \pm 4, \pm 3\rangle$ or
$\langle 0, \pm 5 \rangle$. Notice that type I vectors have an odd
$x$ component and an even $y$ component and vice versa for type II vectors.
\vspace{0.2cm}

Suppose that a primitive pentagon of side length 5 does indeed exist.
Let it have vertices $P_1, P_2,\ldots, P_5$. We must have
$$
\overrightarrow{V} := \overrightarrow{P_1 P_2} + \overrightarrow{P_2 P_3} + \cdots +
\overrightarrow{P_5 P_1} = \overrightarrow{0}
$$
Suppose that among $\overrightarrow{P_1 P_2}, \overrightarrow{P_2 P_3},
\ldots, \overrightarrow{P_5 P_1}$, there are $n$ vectors of type I, and
therefore $(5-n)$ vectors of type II. Because
$5-n \equiv 1 + n \ (\text{mod } 2)$, $n$ and $(5-n)$ have different
parities. If there is an odd number of type I vectors
then $\overrightarrow{V}$ has an odd $x$ component, and if there is an odd
number of type II vectors then $\overrightarrow{V}$ has an odd $y$
component. In both cases it is the case that $\overrightarrow{V}$ cannot be
the zero vector, contradiction. $\blacksquare$\\

\textit{Comments: Frankly a very nice problem, finding the invariant that
one of the components of the vector $\overrightarrow{P_1 P_2} +
\overrightarrow{P_2 P_3} + \cdots + \overrightarrow{P_5 P_1}$ must be odd
allows us to finish the contradiction nicely.}
\end{solution}
\newpage


\begin{problem}{Problem 7}
A \textit{quirky} sequence is a sequence of positive integers
$a_1, a_2, \ldots$ satisfying the following properties:
\begin{itemize}
\item $a_1 = 1$;
\item $a_{i+1} > a_i + 1$ for every $i\geq 1$, and
\item $a_i \pm a_j$ is composite for every $i,j$ with $i > j$.
\end{itemize}
(a) Construct a quirky sequence.
\vspace{0.2cm}

(b) Given a set of numbers $a_1, a_2, \ldots, a_n$ satisfying the quirkyness
properties, can this always be extended into an infinite quirky sequence?
\end{problem}

\begin{solution}{Solution}
(a) Define $a_k = 5^{k-1}$. The first two conditions are clearly satisfied.
Since the two smallest numbers in the sequence
are $a_1 = 1, a_2 = 5$ we have
\[
\begin{aligned}
a_i + a_j &\geq 1 + 5 = 6\\
a_i - a_j &\leq 1 - 5 = -4
\end{aligned}
\]
Notice that $a_k$ consists of only odd terms.
Now since $a_i \pm a_j$ is the sum of two odd integers, it is even. This
coupled with the fact that $|a_i \pm a_j| \geq 4$ implies that the quantity
$a_i \pm a_j$ is in fact composite. $\square$\\

(b) Suppose $a_1, a_2, \ldots, a_n$ is a quirky sequence. Let
$c_1, c_2, \ldots, c_{2n}$ be composite positive integers such
that they are pairwise coprime. If we can find a solution
$a_{n+1} > a_n + 1$ to the system
\[
\begin{aligned}
a_{n+1} + a_1 &\equiv 0 \ (\text{mod } c_1)\\
&\vdots\\
a_{n+1} + a_n &\equiv 0 \ (\text{mod } c_n)\\
a_{n+1} - a_1 &\equiv 0 \ (\text{mod } c_{n+1})\\
&\vdots\\
a_{n+1} - a_n &\equiv 0 \ (\text{mod } c_{2n})\\
\end{aligned}
\]
then the sequence $a_1, a_2, \ldots, a_{n+1}$ is also quirky. However such
a solution is guarenteed due to the Chinese remainder theorem. Now by
iteratively repeating this arguement we can keep on introducing new terms
into the sequence, and so indeed we can create an infinite quirky sequence.
$\blacksquare$
\end{solution}
\newpage

\begin{problem}{Problem 8}
$ABCDEF$ is a cyclic hexagon in which $AF$ and $CD$ are parallel. The lines
$BC$ and $EF$ intersect at $L$, and the lines $AB$ and $DE$ intersect at
$M$. \vspace{0.2cm}

Prove that $LM$ is parallel to $CD$.
\end{problem}

\begin{solution}{Solution}
The arguement splits into two parts.
The following diagram gives an overview of the first part of our arguement.
\begin{center}
\begin{asy}
/* Geogebra to Asymptote conversion, documentation at artofproblemsolving.com/Wiki go to User:Azjps/geogebra */
import graph; size(270);
real labelscalefactor = 0.5; /* changes label-to-point distance */
pen dps = linewidth(0.7) + fontsize(10); defaultpen(dps); /* default pen style */
pen dotstyle = black; /* point style */
real xmin = -7.145383243355196, xmax = 11.636811683872974, ymin = -6.349095835703575, ymax = 12.011874323671082;  /* image dimensions */
pen qqttcc = rgb(0.,0.2,0.8); pen qqzzqq = rgb(0.,0.6,0.); pen ffxfqq = rgb(1.,0.4980392156862745,0.);
pair O = (0.,0.), A = (0.,6.), F = (3.678895565319094,4.73980246629282), B = (-2.4736630207026806,5.466350817502303), C = (-5.986698741850665,0.3992971003196306), D = (4.974123632197761,-3.3553083452362094), L = (1.1397665965708184,10.678208432203846), M = (9.079013303114511,7.95863704415021), G = (2.8748780137761,6.620204243258815);

filldraw(circle(O, 6.), qqttcc + opacity(0.10000000149011612),  qqttcc);
filldraw(arc(C,1.1129951530935935,-18.90878251891273,24.183304184973036)--(-5.986698741850665,0.3992971003196306)--cycle, red + opacity(0.10000000149011612),  red);
filldraw(arc(F,1.1129951530935935,161.09121748108728,204.18330418497305)--(3.678895565319094,4.73980246629282)--cycle, red + opacity(0.10000000149011612),  red);
filldraw(arc(B,1.1129951530935935,235.26606848443078,372.17398178054503)--(-2.4736630207026806,5.466350817502303)--cycle, blue + opacity(0.10000000149011612),  blue);
filldraw(arc((5.9684156499251735,-0.6148289434698667),1.1129951530935935,113.150463602668,250.05837689878223)--(5.9684156499251735,-0.6148289434698667)--cycle, blue + opacity(0.10000000149011612),  blue);
filldraw(arc((5.9684156499251735,-0.6148289434698667),1.1129951530935935,430.0583768987822,473.150463602668)--(5.9684156499251735,-0.6148289434698667)--cycle, red + opacity(0.10000000149011612),  red);
filldraw(arc(B,1.1129951530935935,12.173981780545025,55.266068484430804)--(-2.4736630207026806,5.466350817502303)--cycle, red + opacity(0.10000000149011612),  red);
filldraw(arc(G,1.1129951530935935,113.150463602668,192.173981780545)--(2.8748780137761,6.620204243258815)--cycle, qqzzqq + opacity(0.10000000149011612),  qqzzqq);
filldraw(arc(G,1.1129951530935935,293.15046360266797,372.17398178054503)--(2.8748780137761,6.620204243258815)--cycle, qqzzqq + opacity(0.10000000149011612),  qqzzqq);
filldraw(arc(L,1.1129951530935935,-124.73393151556921,-66.84953639733202)--(1.1397665965708184,10.678208432203846)--cycle, ffxfqq + opacity(0.10000000149011612),  ffxfqq);
filldraw(arc(M,1.1129951530935935,-167.826018219455,-109.94162310121779)--(9.079013303114511,7.95863704415021)--cycle, ffxfqq + opacity(0.10000000149011612),  ffxfqq);
/* draw figures */
draw((xmin, 1.442357584547096*xmin + 9.034257437026493)--(xmax, 1.442357584547096*xmax + 9.034257437026493)); /* line */
draw((xmin, -2.3387571245892658*xmin + 13.343845680302707)--(xmax, -2.3387571245892658*xmax + 13.343845680302707)); /* line */
draw((xmin, 0.21573236856898417*xmin + 6.)--(xmax, 0.21573236856898417*xmax + 6.)); /* line */
draw((xmin, 2.756211809916845*xmin-17.065046644286152)--(xmax, 2.756211809916845*xmax-17.065046644286152)); /* line */
draw(L--M,  linetype("4 4") + red);
draw(A--F,  linetype("4 4") + red);
draw((2.0149366401168147,5.3097879134192825)--(1.7673117989869938,5.1593146041976885),  linetype("4 4") + red);
draw((2.0149366401168147,5.3097879134192825)--(1.9115837663321016,5.580487862095132),  linetype("4 4") + red);
draw(C--D,  linetype("4 4") + red);
draw((-0.3307986973691843,-1.5381189421854184)--(-0.5784235384990053,-1.6885922514070106),  linetype("4 4") + red);
draw((-0.3307986973691843,-1.5381189421854184)--(-0.43415157115389885,-1.2674189935095705),  linetype("4 4") + red);
draw(C--F,  linetype("4 4"));
/* dots and labels */
dot(O, dotstyle);
label("$O$", (0.10823441570224968,0.18600119968622317), NE * labelscalefactor);
dot(A,dotstyle);
label("$A$", (0.10823441570224968,6.245641477640231), NE * labelscalefactor);
dot(F,dotstyle);
label("$F$", (3.7687518080989575,4.984246970800825), NE * labelscalefactor);
dot(B,dotstyle);
label("$B$", (-2.3650881467279583,5.7015105139055855), NE * labelscalefactor);
dot(C,dotstyle);
label("$C$", (-5.877206185378854,0.6559324865479625), NE * labelscalefactor);
dot(D, dotstyle);
label("$D$", (5.079612766186968,-3.1529842595945565), NE * labelscalefactor);
dot((5.9684156499251735,-0.6148289434698667),dotstyle);
label("$E$", (6.068941791159051,-0.35812976404842245), NE * labelscalefactor);
//label("$i$", (-3.3296839460757393,13.418276908687831), NE * labelscalefactor);
dot(L, dotstyle);
label("$L$", (1.2459627944201455,10.870754669384718), NE * labelscalefactor);
dot(M, dotstyle);
label("$M$", (9.185328219821113,8.150099850711491), NE * labelscalefactor);
dot(G, dotstyle);
label("$G$", (2.977288588121291,6.814505666999179), NE * labelscalefactor);
clip((xmin,ymin)--(xmin,ymax)--(xmax,ymax)--(xmax,ymin)--cycle);
/* end of picture */
\end{asy}
\end{center}
%\definecolor{darkgreen}{RGB}{0,133,0}

\textit{Claim 1: $\triangle BGL \sim \triangle EGM$.}\vspace{0.2cm}

\textit{Proof: } Owing to the alternate angle theorem,
${\angle AFC} = {\angle FCD} :=
{\theta}$. Now since opposite angles in a cyclic
quadrilateral sum to $180^\circ$, we have
${\angle FBC} = {\angle FED} =
180 - {\theta}$. But now since angles on a straight line sum
to $180^\circ$ we deduce
${\angle LBG} = {\angle GEM} =
{\theta}$. We also have
${\angle BGL} = {\angle EGM}$ due to
vertically opposite angles being equal.
Now since $\triangle BGL$ and $\triangle EGM$ have two equal angles, their
third angles' must be equal also and so $\triangle BGL \sim \triangle EGM$
by AAA criterion. $\square$ \\

\textit{Claim 2: Quadrilateral $LMEB$ is cyclic.} \vspace{0.2cm}

\textit{Proof: } We note that $\angle LBM = \angle LEM$, and so the line
segment $LM$ subtends two equal angles at points $B$ and $E$.
By the converse
of the "Angles subtending the same segment are equal" theorem, we can deduce
that $LMEB$ are concyclic. $\square$ \\

The following diagram gives an overview of the second part of our
arguement.

\begin{center}
\begin{asy}
/* Geogebra to Asymptote conversion, documentation at artofproblemsolving.com/Wiki go to User:Azjps/geogebra */
import graph; size(270);
real labelscalefactor = 0.5; /* changes label-to-point distance */
pen dps = linewidth(0.7) + fontsize(10); defaultpen(dps); /* default pen style */
pen dotstyle = black; /* point style */
real xmin = -7.389137561570594, xmax = 12.733489266356116, ymin = -7.157870372699978, ymax = 16.24176461576778;  /* image dimensions */
pen ttttff = rgb(0.2,0.2,1.); pen qqwwzz = rgb(0.,0.4,0.6); pen ffzzff = rgb(1.,0.6,1.); pen qqccqq = rgb(0.,0.8,0.); pen ffzzcc = rgb(1.,0.6,0.8);
pair O = (0.,0.), A = (-1.3311258534328012,5.850478951532327), F = (3.1232172370599534,5.123037584298166), C = (-5.933347405288747,-0.8918455965879346), D = (5.3413516262621705,-2.7331232691970677), B = (-4.573826133661758,3.883312310262045), G = (1.6866967014810506,7.681224733217836), L = (-1.7500564535895504,13.801469823739785), M = (8.906580392780866,12.061128246956558);

filldraw(circle(O, 6.), ttttff + opacity(0.05000000074505806),  ttttff);
filldraw(circle((2.6407021743535424,7.190337379924182), 7.936361478569885), qqwwzz + opacity(0.10000000149011612),  qqwwzz);
filldraw(arc(A,1.9932448981981477,350.7248714177551,391.24286847748385)--(-1.3311258534328012,5.850478951532327)--cycle, ffzzff + opacity(0.10000000149011612),  ffzzff);
filldraw(arc(A,0.8542478135134919,-148.75713152251615,-9.275128582244916)--(-1.3311258534328012,5.850478951532327)--cycle, qqccqq + opacity(0.10000000149011612),  qqccqq);
filldraw(arc((6.,0.),1.4237463558558199,479.3158869886146,519.8338840483434)--(6.,0.)--cycle, ffzzcc + opacity(0.10000000149011612),  ffzzcc);
filldraw(arc(M,1.4237463558558199,170.72487141775508,211.24286847748385)--(8.906580392780866,12.061128246956558)--cycle, ffzzcc + opacity(0.10000000149011612),  ffzzcc);
/* draw figures */
draw(A--F,  linetype("4 4") + red);
draw((1.1068154878902132,5.452337329627645)--(0.8547405658684547,5.2338345126232815),  linetype("4 4") + red);
draw((1.1068154878902132,5.452337329627645)--(0.9373508177586982,5.7396820232072105),  linetype("4 4") + red);
draw(C--D,  linetype("4 4") + red);
draw((-0.08522809343664976,-1.846905371180105)--(-0.3373030154584083,-2.065408188184467),  linetype("4 4") + red);
draw((-0.08522809343664976,-1.846905371180105)--(-0.2546927635681648,-1.559560677600538),  linetype("4 4") + red);
draw((xmin, 3.5123818997957805*xmin + 19.948336434948516)--(xmax, 3.5123818997957805*xmax + 19.948336434948516)); /* line */
draw((xmin, -1.7808218438650776*xmin + 10.684931063190465)--(xmax, -1.7808218438650776*xmax + 10.684931063190465)); /* line */
draw((xmin, 0.6066446082804134*xmin + 6.65799927346)--(xmax, 0.6066446082804134*xmax + 6.65799927346)); /* line */
draw((xmin, 4.14959389284846*xmin-24.897563357090768)--(xmax, 4.14959389284846*xmax-24.897563357090768)); /* line */
draw(F--C,  linetype("4 4"));
draw(L--M,  linetype("4 4") + red);
draw(B--(6.,0.));
/* dots and labels */
dot(O, dotstyle);
label("$O$", (0.10122279014072529,0.21375819973222002), NE * labelscalefactor);
dot((6.,0.),dotstyle);
label("$E$", (6.109432411852285,0.27070805396645264), NE * labelscalefactor);
dot(A,dotstyle);
label("$A$", (-1.208623857246629,6.136543040092414), NE * labelscalefactor);
dot(F,dotstyle);
label("$F$", (3.233464773023529,5.39619493504739), NE * labelscalefactor);
dot(C,dotstyle);
label("$C$", (-5.821562050219486,-0.6120146866641533), NE * labelscalefactor);
dot(D, dotstyle);
label("$D$", (4.913485472933396,-2.4913598763938305), NE * labelscalefactor);
dot(B,dotstyle);
label("$B$", (-4.454765548597898,4.171773069011389), NE * labelscalefactor);
dot(G, dotstyle);
label("$G$", (1.8097184171677092,7.901988521353626), NE * labelscalefactor);
dot(L, dotstyle);
label("$L$", (-1.6357477640033748,14.024097851533636), NE * labelscalefactor);
dot(M, dotstyle);
label("$M$", (9.013874977798158,12.28712729738954), NE * labelscalefactor);
label("$\theta$", (0.7561461138344024,6.079593185858182), NE * labelscalefactor,ffzzff);
label("$180-\theta$", (-1.6357477640033748,4.541947121533901), NE * labelscalefactor,qqccqq);
label("$\theta$", (4.457886639059534,0.9825812318943606), NE * labelscalefactor,ffzzcc);
label("$\theta$", (7.049105006717126,11.546779192344516), NE * labelscalefactor,ffzzcc);
clip((xmin,ymin)--(xmin,ymax)--(xmax,ymax)--(xmax,ymin)--cycle);
/* end of picture */
\end{asy}
\end{center}

Let $\angle GAF = \theta$. Since $B, A$ and $G$ lie on a straight line,
$\angle BAF = 180 - \theta$. \vspace{0.2cm}

Now consider the cyclic quadrilateral
$ABEF$. Since opposite angles in a cyclic quadrilateral sum to $180^\circ$
we have $\angle BEL = \theta$. \vspace{0.2cm}

Now look at the circle through $LMEB$. The segment $LB$ subtends the points
$M$ and $E$, so by the "Angles subtending the same segment are equal"
theorem, $\angle BEL = \angle LMB = \theta$. \vspace{0.2cm}

Now finally by the converse of the "Alternate angles are equal" theorem
it must be the case that $AF \parallel LM$. $\blacksquare$
\end{solution}

\end{document}
